% ============================================================
%  2.3  Hardware Acceleration for Deep Learning
% ============================================================
\section{Hardware Acceleration for Deep Learning}
\label{sec:bg-hw-accel}

The operations central to DNNs---matrix multiplications and
convolutions---are, at their core, simple linear-algebra primitives.
However, they involve large matrices and tensors and are executed
hundreds or thousands of times within a single network, leading to
exceptionally high computational and memory demands.
Conversely, the simplicity of these kernels is an asset: they are
\emph{uniform}, exhibit \emph{regular} data dependencies and access
patterns, and offer \emph{abundant parallelism}.
This regularity opens numerous opportunities for execution speedups
and energy-saving optimizations through specialized hardware.


% ------------------------------------------------------------
\subsection{From CPUs to Spatial Architectures}
\label{sec:cpu-gpu-sa}
% ------------------------------------------------------------

General-purpose CPUs can execute DNN workloads through their SIMD
vector units, but they are fundamentally constrained by a control-heavy
micro-architecture designed for diverse, irregular code.
GPUs improve on CPUs by providing thousands of lightweight threads
operating in a Single-Instruction, Multiple-Thread (SIMT) model,
achieving higher throughput on data-parallel workloads.
However, GPUs still rely on general-purpose cores and a cache-based
memory hierarchy that does not explicitly exploit the reuse patterns
specific to DNN computations.

\textbf{Spatial Architectures} (SAs) represent a fundamentally different
approach: instead of relying on shared, general-purpose cores, they
employ a multidimensional array of dedicated \emph{Processing Elements}
(PEs) paired with an explicit on-chip memory hierarchy
(registers, scratchpads, FIFOs) engineered to exploit the regular
data dependencies and access patterns inherent to DNN kernels.
Each PE is a simple functional unit executing \emph{multiply-accumulate}
(MAC) operations, and the architecture is designed such that data is
passed between PEs and memories through controlled, predictable
interconnects, maximizing \emph{data reuse} and minimizing costly
accesses to off-chip DRAM.

This specialized design philosophy has led to substantial improvements
in both performance and energy efficiency.
For example, Google's TPUv1~\cite{tpuv1} achieved a
$15$--$30\times$ increase in computational throughput and a
$30$--$80\times$ improvement in energy efficiency compared to
contemporary CPUs and GPUs executing the same inference workloads.


% ------------------------------------------------------------
\subsection{Abstracting Spatial Architectures}
\label{sec:sa-abstraction}
% ------------------------------------------------------------

Despite the diversity of existing spatial architectures, their structure
can be captured by a unified abstraction consisting of three
\emph{level types} arranged in a hierarchy from outermost (closest to
off-chip memory) to innermost (closest to computation):

\begin{enumerate}
  \item \textbf{Memory levels} store tiles of one or more operands
    (inputs, weights, outputs) and are characterized by their
    \emph{capacity} (in number of entries), \emph{bandwidth} (operands
    per clock cycle), and per-access \emph{energy cost} (read and write,
    in pJ).
    A memory level may \emph{bypass} certain operands, meaning it does
    not store them and delegates their handling to adjacent levels.
    The outermost level of any architecture must be a memory level,
    typically representing DRAM or another form of off-chip storage.

  \item \textbf{Fanout (spatial) levels} model the physical replication
    of all subsequent levels into a multidimensional mesh of instances.
    A fanout level is defined by its \emph{mesh size} (i.e., the number
    of parallel instances it creates) and the set of loop dimensions that
    can be spatially unrolled across those instances.
    Fanout levels are the mechanism through which spatial architectures
    achieve parallelism: each instance in the mesh executes the same
    inner loop nest concurrently, but on different tiles of data.

    Fanout levels additionally specify hardware support for
    \emph{spatial multicast} (one read from an outer memory feeds all
    instances requiring the same data) and \emph{spatial reduction}
    (partial results from multiple instances are accumulated into a
    single write to an outer memory).
    When multicast or reduction support is absent, each instance must
    independently fetch its own copy of shared data, significantly
    increasing bandwidth pressure on outer memories.

  \item \textbf{Compute levels} represent the functional units (PEs)
    that perform the MAC operations.
    A compute level is defined by its \emph{energy per MAC} and the
    number of \emph{clock cycles per MAC}.
    The compute level must be the innermost level in the hierarchy and
    appears exactly once.
\end{enumerate}

Fanout and compute levels are collectively referred to as
\textbf{spatial levels}, because they create multiple concurrent
instances of the computation.
The key insight of this abstraction is that any level---whether memory
or spatial---can be described through the same lens: each level holds
a copy of the loop nest with its own \emph{factors} (the number of
iterations unfolding at that level for each dimension) and its own
\emph{dataflow} (the order of those loops).
Combined, the factors and dataflows of all levels determine the
complete mapping (see Section~\ref{sec:bg-mapping}).


% ------------------------------------------------------------
\subsection{Data Reuse in Spatial Architectures}
\label{sec:data-reuse}
% ------------------------------------------------------------

The primary advantage of spatial architectures over general-purpose
processors lies in their ability to systematically exploit
\emph{data reuse}---the phenomenon whereby a single data element
is accessed multiple times during a computation, and only the first
access need be serviced by an expensive outer memory.
Three forms of reuse are central to the architecture model:

\begin{itemize}
  \item \textbf{Temporal reuse} occurs when a data element resides in a
    memory level and is read or updated multiple times by consecutive
    iterations at that level, without being evicted back to outer memory.
    The classical notion of \emph{stationarity} (weight-stationary,
    input-stationary, output-stationary) is a specific form of temporal
    reuse: an operand is stationary at a given level when the innermost
    loops at that level iterate over dimensions \emph{orthogonal} to
    it~(cf.\ Definition~\ref{def:coupled-orthogonal}).

  \item \textbf{Spatial multicast reuse} occurs across the instances
    of a fanout level.
    When a fanout spatially unrolls a dimension that is orthogonal to
    an operand, all instances require the same tile of that operand.
    With hardware multicast support, a single read from the outer memory
    is broadcast to all instances, amortizing the access cost.

  \item \textbf{Spatial reduction reuse} is the dual of multicast:
    when a fanout spatially unrolls a dimension that is coupled to an
    output (i.e., multiple instances produce partial sums for the same
    output tile), hardware reduction support allows those partial sums
    to be accumulated and written back as a single result.
\end{itemize}

Maximizing reuse is the central objective of the mapping problem:
accessing off-chip DRAM can cost hundreds of times the energy
of a PE-local register read (e.g., $64\;\text{pJ}$ vs.
${\sim}0.1\;\text{pJ}$ per operand~\cite{eyeriss}), making data
movement the dominant contributor to total energy unless reuse is
carefully orchestrated.


% ------------------------------------------------------------
\subsection{Examples of Spatial Architectures}
\label{sec:sa-examples}
% ------------------------------------------------------------

To ground the abstraction, this section briefly describes three
well-known spatial architectures that illustrate different points
in the design space.

\paragraph{Eyeriss~v1~\cite{eyeriss}.}
Eyeriss is a dataflow-flexible CNN accelerator featuring a
$14 \times 12 = 168$ PE array organized in a two-dimensional mesh.
Each PE contains a local register file for inputs, weights, and partial
sums.
A shared Global Buffer (108~KB SRAM) sits between the PE array and
off-chip DRAM, storing activations while weights are streamed directly
from DRAM to the PEs.
Eyeriss supports a Row-Stationary dataflow that maximizes the reuse of
all three operands simultaneously at the PE level, but its flexible
architecture allows the mapper to explore other dataflow choices as well.

\paragraph{Simba~\cite{simba}.}
Simba is a multi-chip-module (MCM) DNN inference accelerator consisting
of multiple chiplets, each containing its own PE array and local SRAM.
Data flows through three levels of the memory hierarchy: DRAM, per-chiplet
SRAM (``PE Buffer''), and per-PE register files, with two fanout levels
(inter-chiplet and intra-chiplet) providing spatial parallelism.
This deeper hierarchy makes Simba's map-space significantly richer than
Eyeriss's, offering more opportunities for data reuse but also a larger
space of sub-optimal mappings.

\paragraph{TPUv1~\cite{tpuv1}.}
Google's first-generation Tensor Processing Unit employs a
$256 \times 256 = 65{,}536$-element systolic array, where weights are
preloaded into registers and activations flow through PEs in a
pipelined, PE-to-PE fashion.
A 24~MB Unified Buffer stores input activations and final outputs,
while a separate 4~MB Weight FIFO stages weight tiles between DRAM and
the array.
The systolic data flow imposes a weight-stationary constraint on the
innermost level, but tiling and loop ordering at the outer memory levels
remain degrees of freedom for the mapper.
\vspace{0.3em}

\noindent
Despite their architectural differences, all three designs share the
same fundamental structure: a hierarchy of memory levels interspersed
with spatial fanout levels leading to a compute level.
This structural uniformity is what allows a single abstract model to
capture them all, and it is the foundation upon which the mapping
framework presented in this thesis is built.
